\documentclass{article}
\usepackage[margin=2cm]{geometry}
\usepackage{amsmath}
\usepackage{graphicx}
\usepackage{booktabs}
\usepackage[utf8]{inputenc}
\usepackage{ragged2e}
\setlength{\emergencystretch}{3em}
\begin{document}
\sloppy

\noindent\hfill\includegraphics[height=1.8em]{assets/figure_header_logo.png}\par\noindent\hrule\vspace{4pt}

e;   e, for all    j. The free energy in the specific case of a vector multiplet in the
adjoint, a single antisymmetric hypermultiplet, and Nf fundamental hypermultiplets then is


\begin{equation}
\begin{array} { c } { F ( \lambda _ { i } ) = \sum _ { i \neq j } \left [ F _ { V } ( \lambda _ { i } - \lambda _ { j } ) + F _ { V } ( \lambda _ { i } + \lambda _ { j } ) + F _ { H } ( \lambda _ { i } - \lambda _ { j } ) + F _ { H } ( \lambda _ { i } + \lambda _ { j } ) \right ] } \\ { \qquad \qquad + \sum _ { i } \left [ F _ { V } ( 2 \lambda _ { i } ) + F _ { V } ( - 2 \lambda _ { i } ) + N _ { f } F _ { H } ( \lambda _ { i } ) + N _ { f } F _ { H } ( - \lambda _ { i } ) \right ] \qquad \qquad ( 1 ) } \end{array}
\end{equation}


The next step is to look for extrema of this function in the specific case of  ; > O for all
. Extrema in the case of non-positive  ; can be obtained from these through action of the
Weylggroup. To calculate the extrema, one first assumes that as N    oo, the vevs scale as
 ; = N x; for   > 0 and x; of order O(N ). One then introduces a density function


\begin{equation}
\rho ( x ) = \frac { 1 } { N } \sum _ { i = 1 } ^ { N } \delta ( x - x _ { i } ) \qquad ( 2 )
\end{equation}


which in the continuum limit should approach an Ll function normalized as


\begin{equation}
\int d x ~ \rho ( x ) = 1 \qquad ( 3 )
\end{equation}


In terms of thiss densityfunction,onefinds that


\begin{equation}
F \approx - \frac { 9 \pi } { 8 } N ^ { 2 + \alpha } \int d x d y ~ \rho ( x ) \rho ( y ) \left ( | x - y | + | x + y | \right ) + \frac { \pi ( 8 - N _ { f } ) } { 3 } N ^ { 1 + 3 \alpha } \int d z ~ \rho ( x ) ~ | x | ^ { 3 } \quad ( 4 )
\end{equation}


where the large argument expansions have been used, and terms subleading in N have been
dropped. This only has non-trivial saddle points when both terms scale the same with N,
which demands that   = 1/2 andgives the famous result that F  x N5/2.EExtremizingthe
gives


\begin{equation}
F \approx - \frac { 9 \sqrt { 2 } \pi N ^ { 5 / 2 } } { 5 \sqrt { 8 - N _ { f } } } \qquad ( 5 )
\end{equation}


This value of the free energy is to be identifiedwith theerenormalizedon-shell action of the
supersymmetric AdSs solution. This identification yields the following relationbetweenthe
six-dimensionalNewton'sconstantG6and theparametersNaand Nf of the dual SCFT,


\begin{equation}
G _ { 6 } = \frac { 5 \pi \sqrt { 8 - N _ { f } } } { 2 7 \sqrt { 2 } } ~ N ^ { - 5 / 2 } \qquad ( 6 )
\end{equation}

\end{document}
